\section{Cross-comparison with SeQUeNCe\label{sequence}}

We compare QuantumSavory against SeQUeNCe on a resource-management example that
is deliberately small enough to inspect, yet complex enough to exercise several
independent layers of a quantum-network simulator.  The example combines
memory-resource allocation, Barrett-Kok elementary-link generation,
entanglement swapping, and BBPSSW purification in a single run.  Agreement on
this example is therefore more informative than agreement on an isolated link:
it checks not only the elementary generation model, but also how generated
resources are reserved, consumed, transformed, and reported.  The implementation
and all configuration files used for this comparison are available in the
companion repository.\footnote{\url{https://lbacciottini.github.io/SeQUeNCe_QSavory_comparison/}}

The scenario is adapted from the SeQUeNCe resource-management
tutorial.\footnote{\url{https://sequence-rtd-tutorial.readthedocs.io/stable/tutorial/chapter4/resource_management.html}}
It consists of a three-node repeater chain, denoted
$r_1-r_2-r_3$, with midpoint Bell-state measurement stations on the two
elementary links.  Two traffic demands are active at the same time.  The first
flow asks for direct $r_1-r_2$ elementary pairs.  The second flow asks for
end-to-end $r_1-r_3$ pairs, which are created by generating one elementary pair
on each adjacent link and then performing entanglement swapping at $r_2$.
Only these swapped $r_1-r_3$ pairs are purified; elementary pairs are never
distilled in this example.  The BBPSSW operation is ideal, with classical
initial-handshake timing included but no local gate noise.  Memories are
multiplexed: each reserved pair of
memory slots may attempt elementary-link generation independently and in
parallel with the other reserved slot pairs.  This is the intended meaning of
memory multiplexing throughout this section.  We do not model memory
decoherence or memory cutoff, so all observed differences come from generation
statistics, protocol timing, and the chosen quantum-state representation.

\begin{table}[t]
    \caption{Shared hardware and protocol configuration used in the comparison.
    The same configuration is adapted into both simulators.  SeQUeNCe uses a
    Bell-diagonal-state (BDS) representation, which is Werner-like for this
    comparison; QuantumSavory is run both with a Werner/depolarized state and
    with its Barrett-Kok-specific state representation.}
    \label{tab:sequence_qsavory_config}
    \begin{ruledtabular}
    \begin{tabular}{ll}
    Quantity & Value \\
    \hline
    Topology & Three nodes $r_1-r_2-r_3$ with two midpoint BSM stations \\
    Link length & Sweep over $5,10,\ldots,50$ km \\
    Signal speed & $2\times 10^5$ km/s \\
    Memories & $20$ at $r_1$, $40$ at $r_2$, $10$ at $r_3$ \\
    Flow reservations & $10$ multiplexed memories per flow endpoint \\
    Fiber attenuation & $0.2$ dB/km \\
    Source efficiencies & emission $1$, collection $0.5$, conversion $1$ \\
    Detector model & efficiency $0.8$, dark-click probability $5\times10^{-4}$ \\
    Detector timing & count rate $25$ MHz, resolution $150$ ps \\
    Barrett-Kok visibility & $0.99$ \\
    Barrett-Kok success model & $p_{\mathrm{BK}}=p_{\mathrm{click}}^2/2$ \\
    Swapping & ideal, with success probability $1$ \\
    Purification & ideal BBPSSW on swapped $r_1-r_3$ pairs, with initial handshake \\
    QuantumSavory states & \texttt{BarrettKokBellPair} and \texttt{DepolarizedBellPair} \\
    SeQUeNCe state representation & BDS/Werner-like representation \\
    \end{tabular}
    \end{ruledtabular}
\end{table}

\subsection{Shared Barrett-Kok model}

Both simulators are driven by the same shared Barrett-Kok generation model:
the elementary success probability, effective attempt time, and raw-state
fidelity are derived once from the common hardware configuration and then
adapted into each simulator.  The exact density matrix used by QuantumSavory's
\texttt{BarrettKokBellPair}, including the definitions of every optical
efficiency and detector parameter, is derived in the companion
documentation.\footnote{\url{https://lbacciottini.github.io/SeQUeNCe_QSavory_comparison/0.1.0/explanation/barrett-kok-state/}}
In addition to this exact model, we instantiate a Werner-state model in
QuantumSavory by using \texttt{DepolarizedBellPair} with the same raw fidelity
as the exact Barrett-Kok pair.  This second model is included because Werner
states are a common approximation in the quantum-networking literature and
because SeQUeNCe is most naturally run in a Bell-diagonal-state representation,
which is Werner-like for this comparison.  Adding the same abstraction to
QuantumSavory gives a direct representation-level validation between the two
simulators, while the exact run separately shows the effect of preserving the
full Barrett-Kok density-matrix structure.

\subsection{QuantumSavory implementation}

The QuantumSavory implementation is most naturally expressed as a small set of
asynchronous protocol processes.  We create one register per network node and
connect them with a \texttt{RegisterNet} whose classical delay is derived from
the speed-of-light delay in the shared configuration.  We then start an
\texttt{EntanglementTracker} on every node so that the metadata describing the
remote endpoint of an entangled pair is updated automatically after swaps.

Elementary generation is multiplexed over the reserved memories.  Concretely,
for each paired set of reserved memory slots, the implementation starts an
independent fixed-slot \texttt{EntanglerProt}.  Each process samples successful
Barrett-Kok generation with the shared success probability; its attempt time
is the SeQUeNCe-equivalent effective attempt duration, including the fact that
the second optical round is only reached when the first round succeeds.  The
production implementation uses the following pattern.
\begin{lstlisting}[language=Julia]
function install_multiplexed_entanglers!(
    sim, net, nodeA, nodeB, slotsA, slotsB,
    success_prob, attempt_time, pairstate,
)
    for (slotA, slotB) in zip(slotsA, slotsB)
        @process EntanglerProt(
            sim, net, nodeA, nodeB;
            chooseslotA=slotA,
            chooseslotB=slotB,
            success_prob=success_prob,
            attempt_time=attempt_time,
            pairstate=pairstate,
            retry_lock_time=nothing,
        )()
    end
end
\end{lstlisting}
The keyword arguments \texttt{chooseslotA} and \texttt{chooseslotB} are the
important part of the construction: they bind each entangling process to one
reserved pair of memories, so different reserved memories may attempt
generation concurrently.

The raw-state model is orthogonal to the protocol logic.  The same entanglers,
swapper, and purifier are used whether the elementary pair is represented by
the exact Barrett-Kok state or by the Werner approximation:
\begin{lstlisting}[language=Julia]
function raw_pair_state(resolved, model)
    if model == "exact"
        return BarrettKokBellPair(
            resolved["derived"]["arm_transmissivity"],
            resolved["derived"]["arm_transmissivity"],
            resolved["detectors"]["dark_count_probability"],
            resolved["detectors"]["efficiency"],
            resolved["barrett_kok"]["mode_matching_visibility"],
            resolved["barrett_kok"]["parity_bit"],
        )
    end

    return DepolarizedBellPair(;
        F=resolved["derived"]["barrett_kok_raw_fidelity"],
    )
end
\end{lstlisting}
This is the intended role of \texttt{StatesZoo}: the protocol stack consumes an
abstract two-qubit state, while the physical model of that state can be changed
without rewriting the control flow.

After elementary generation, a single \texttt{SwapperProt} runs at the middle
node $r_2$.  It is restricted to the memory ranges reserved for the end-to-end
flow and performs ideal swapping.  Finally, purification is installed only
between the end nodes:
\begin{lstlisting}[language=Julia]
@process SwapperProt(
    sim, net, 2;
    nodeL=1, nodeH=3,
    chooseslots=swap_slots,
    retry_lock_time=nothing,
    local_busy_time=swap_busy,
)()

@process BBPSSWProt(
    sim, net, 1, 3;
    retry_lock_time=nothing,
    initial_handshake=true,
)()
\end{lstlisting}
The \texttt{initial\_handshake=true} option makes the BBPSSW timing include the
same initial classical coordination step as in the SeQUeNCe implementation.
The \texttt{retry\_lock\_time=nothing} choices make the swapper and purifier
event-driven: when no valid pair is available, they wait for the relevant
metadata to change instead of polling at a fixed interval.

\subsection{SeQUeNCe implementation}

The SeQUeNCe implementation realizes the same shared configuration at a lower
level of detail.  It instantiates quantum routers, midpoint BSM nodes,
classical and quantum channels, resource-manager rules, Barrett-Kok endpoint
protocols, entanglement-swapping protocols, and BBPSSW purification rules.  The
same memory reservations determine which memories may participate in each flow,
and purification rules are restricted to swapped $r_1-r_3$ pairs.  SeQUeNCe is
run in its Bell-diagonal-state representation, matching the Werner-like
abstraction used by the \texttt{DepolarizedBellPair} QuantumSavory run.

The main difference between the implementations is the level at which the event
trace is exposed.  SeQUeNCe naturally exposes a detailed trace of resource
manager requests, protocol negotiations, optical emissions, BSM outcomes, and
state updates.  The QuantumSavory implementation above exposes a higher-level
trace in which the same behavior is compressed into asynchronous protocol
processes and state objects.  This is a modeling choice, not a limitation:
QuantumSavory could be instrumented to emit a lower-level event trace, and
SeQUeNCe could be wrapped to present a higher-level trace.  The two
abstraction levels used here are simply the most idiomatic ones for the
respective simulators.

\subsection{Comparison results}

\begin{figure*}[t]
    \centering
    \begin{minipage}{0.48\textwidth}
        \centering
        \textbf{(a)}\\[0.5ex]
        \includegraphics[width=\linewidth]{completion_time_by_link_length.pdf}
    \end{minipage}
    \hfill
    \begin{minipage}{0.48\textwidth}
        \centering
        \textbf{(b)}\\[0.5ex]
        \includegraphics[width=\linewidth]{average_fidelity_by_link_length.pdf}
    \end{minipage}
    \caption{Cross-comparison between SeQUeNCe and QuantumSavory over a
    link-length sweep.  (a) Mean completion time for the requested end-to-end
    resources.  (b) Mean fidelity of distilled end-to-end pairs.  Error bars
    show confidence intervals over seeded repetitions.}
    \label{fig:sequence_qsavory_comparison}
\end{figure*}

Figure~\ref{fig:sequence_qsavory_comparison}(a) shows that the completion times
of the three simulator series agree within confidence intervals over the full
link-length sweep.  This is the main timing check: replacing SeQUeNCe's
lower-level Barrett-Kok event trace by the higher-level QuantumSavory
\texttt{EntanglerProt} abstraction does not introduce a statistically
significant timing discrepancy for this experiment, provided that the
compressed attempt time is derived from the same handshake and propagation
delays.

The fidelity comparison in Fig.~\ref{fig:sequence_qsavory_comparison}(b)
separates two effects.  SeQUeNCe and the QuantumSavory Werner run agree closely
because both use a Bell-diagonal or Werner-like representation for the
elementary pair state.  By contrast, the QuantumSavory exact run produces a
visibly different distilled-pair fidelity because \texttt{BarrettKokBellPair}
retains the physical structure of the Barrett-Kok density matrix rather than
compressing it to a single Werner parameter.  In this example, therefore, the
choice of protocol abstraction has no significant effect on the reported
timing metrics, while the choice of quantum-state representation has a much
larger effect on the predicted distilled-pair fidelity.
